\section{A finite-size criterion for the cube condition}\label{app:finite}

This appendix proves the finite-size criterion used in
Sec.~\ref{sec:massgen}: a gap for the dynamics that resamples the blocks of
one patch of $m^4$ blocks, uniform in the configuration outside the patch,
gives the cube condition on every large torus. Blocks that share no
plaquette give commuting projections, and this replaces the disjointness
used in the finite-size criteria of Knabe~\cite{Kna88} and of Gosset and
Mozgunov~\cite{GM16} for quantum spin systems.\footnote{In those criteria
the local terms of a Hamiltonian act on overlapping sets of sites, and the
gap of the whole system is bounded below by the gap of a finite patch minus
a correction that decays with the size of the patch. Here the local terms
are the complements $I-\mathcal P_j$ of conditional expectations, and two
of them commute as soon as the blocks are at distance at least two.}

\emph{Setting.} The notation is that of Sec.~\ref{sec:cube} and of the
setting of Sec.~\ref{sec:massgen}. Fix integers $d\ge2$ and $m\ge4$ with
$d\mid L$, put $K=L/d$, $n=2d-1$, and fix a shift $a\in[0,d)^4$. For $j$ in
the coarse torus $T_K=(\Z/K\Z)^4$ let $r_j=a+dj$, let $B_j=V_{r_j}$ be the
set of links with base point in $r_j+[0,n)^4$, and write $\mathcal P_j=
\mathcal P_{r_j}$. Displacements $v=j'-j$ in $T_K$ are represented in
$(-K/2,K/2]^4$. Two sets $A,A'$ of links are \emph{non-interacting} if no
plaquette contains a link of each; for a set $A$ of links, $P_A$ is the
conditional expectation under $\mu$ given the links outside $A$.

\begin{lemma}\label{lem:fcommute}
\begin{enumerate}[label=\textup{(\roman*)},leftmargin=*]
\item For non-interacting sets $A,A'$ of links, $P_AP_{A'}=P_{A'}P_A=
P_{A\cup A'}$, and $(I-P_A)(I-P_{A'})$ is an orthogonal projection.
\item If $K\ge4$ and $|j'-j|_\infty\ge2$, then $B_j$ and $B_{j'}$ are
non-interacting.
\item The blocks $B_j$, $j\in T_K$, cover $E$.
\item If $f\in L^2(\mu)$ and $\mathcal P_jf=f$ for every $j\in T_K$, then
$f$ is almost surely constant.
\end{enumerate}
\end{lemma}

\begin{proof}
(i) Every link lies on a plaquette, so $A\cap A'=\emptyset$. Given the
links outside $A\cup A'$, the law of $(U_A,U_{A'})$ has density
proportional to $\prod_{p\cap(A\cup A')\ne\emptyset}\varphi(U_p)$, and each
factor involves the links of exactly one of $A$, $A'$. The density is
therefore a product of a function of $U_A$ and a function of $U_{A'}$,
with coefficients depending on the links outside $A\cup A'$. Hence $U_A$
and $U_{A'}$ are conditionally independent, the conditional law
$\gamma_A$ of $U_A$ given the links outside $A$ does not depend on
$U_{A'}$, and likewise for $A'$. So $P_AP_{A'}f$ and $P_{A'}P_Af$ are both
the integral of $f$ against $\gamma_A\otimes\gamma_{A'}$, which is
$P_{A\cup A'}f$, and $(I-P_A)(I-P_{A'})=I-P_A-P_{A'}+P_{A\cup A'}$ is
self-adjoint and idempotent.

(ii) Let $|v_i|\ge2$ for $v=j'-j$. In the coordinate $i$ the base points of
the links of $B_j$ and of $B_{j'}$ lie in $dj_i+a_i+[0,2d-2]$ and in
$dj_i+a_i+dv_i+[0,2d-2]$ modulo $L$. If $v_i\ge2$, their distance in
$\Z/L\Z$ is at least $\min\{(v_i-2)d+2,(K-v_i-2)d+2\}\ge2$, since $v_i\le
K/2$ and $K\ge4$; the case $v_i\le-2$ is symmetric. A plaquette meeting both
blocks would contain two links whose base points differ by at least two in
the coordinate $i$, while the base points of the links of a plaquette
differ by at most one in each coordinate.

(iii) The sets $r_j+[0,d)^4$ partition $\Lambda_L$, and $[0,d)\subset
[0,n)$.

(iv) The density of $\mu$ with respect to the Haar measure $h$ of $G^E$ is
bounded above and below by positive constants, so $\mu$ and $h$ have the
same null sets and $L^2(\mu)=L^2(h)$. If $\mathcal P_jf=f$, then $f$ is
almost surely a function of $U_{E\setminus B_j}$, and it is fixed by the
conditional expectation $Q_j$ of $h$ given $U_{E\setminus B_j}$, which
integrates out $U_{B_j}$. By Fubini's theorem the $Q_j$ commute and their
product integrates out the links of $\bigcup_jB_j=E$, so $f=\int f\,dh$
almost surely.
\end{proof}

In four dimensions the displacements $v$ with $|v|_\infty=1$ are $80$ in
number, and those with $|v|_\infty\ge2$ give commuting projections. For
$s\in T_K$ let $S_s=s+[0,m-1]^4$ be the \emph{patch} at $s$ and $S^\circ_s=
s+[0,m-2]^4$ the set of its $(m-1)^4$ block indices, and put
\[
H=\sum_{j\in T_K}(I-\mathcal P_j),\qquad H_s=\sum_{j\in S^\circ_s}(I-
\mathcal P_j).
\]
These are bounded nonnegative operators on $L^2(\mu)$ and preserve the
subspace $\mathcal A$ of gauge-invariant functions.

\begin{lemma}\label{lem:knabe}
Let $m\ge4$, $p=m-1$, $K\ge2m$, let $\mathcal K$ be $L^2(\mu)$ or
$\mathcal A$, and suppose that $H_s^2\ge\gamma H_s$ on $\mathcal K$ for
every $s\in T_K$, with $\gamma>0$. Then $H^2\ge\varepsilon H$ on
$\mathcal K$ with
\[
\varepsilon=1-\frac{(1-\gamma)p^4+72p^3-168p^2+96p-16}{p^3(p-2)}\ge
\frac{\gamma p-74}{p-2}.
\]
In particular $\varepsilon\ge\gamma/2$ if $m-1\ge148/\gamma$.
\end{lemma}

\begin{proof}
All operators below preserve $\mathcal K$, and a restriction of an
orthogonal projection to an invariant closed subspace is an orthogonal
projection; so it suffices to argue on $L^2(\mu)$ with inequalities read on
$\mathcal K$. Write $\mathbf P_j=I-\mathcal P_j$ and $A_v=\sum_j\mathbf P_j
\mathbf P_{j+v}$. Since $K\ge2m$, two indices $j$ and $j+v$ lie in a common
$S^\circ_s$ only if $|v_i|\le m-2$ for all $i$, and then for exactly
$c_v=\prod_{i=1}^4(p-|v_i|)$ points $s$. Let $O$ be the set of the $80$
displacements with $|v|_\infty=1$, $D$ the set of those with
$|v|_\infty\ge2$, and $D'$ the set of $v\in D$ with all $|v_i|\le m-2$. For
$v\in D$, $\mathbf P_j\mathbf P_{j+v}$ is an orthogonal projection by
Lemma~\ref{lem:fcommute}(i),(ii), so $A_v\ge0$. Summing the hypothesis over
$s$,
\[
\sum_sH_s^2=c_0H+\sum_{v\in O}c_vA_v+\sum_{v\in D'}c_vA_v\ge\gamma\sum_sH_s
=\gamma c_0H .
\]
Put $c_*=p^3(p-2)$; then $c_v\le c_*$ for $v\in D'$, since one factor of
$c_v$ is at most $p-2$ and the others are at most $p$. For orthogonal
projections $\mathbf P,\mathbf P'$, $(\mathbf P\pm\mathbf P')^2\ge0$ gives
$-(\mathbf P+\mathbf P')\le\mathbf P\mathbf P'+\mathbf P'\mathbf P\le
\mathbf P+\mathbf P'$, so $-2H\le A_v+A_{-v}\le2H$. As $O$ is symmetric and
$c_{-v}=c_v$,
\[
\sum_{v\in O}(c_*-c_v)A_v\ge-\Sigma H,\qquad\Sigma=\sum_{v\in O}|c_*-c_v|,
\]
and therefore
\[
c_*\sum_{v\in O}A_v=\sum_{v\in O}c_vA_v+\sum_{v\in O}(c_*-c_v)A_v\ge
-\bigl((1-\gamma)c_0+\Sigma\bigr)H-\sum_{v\in D'}c_vA_v .
\]
Since $H^2=\sum_vA_v=H+\sum_{v\in O}A_v+\sum_{v\in D}A_v$,
\[
H^2\ge\Bigl(1-\frac{(1-\gamma)c_0+\Sigma}{c_*}\Bigr)H+\sum_{v\in D'}\Bigl(1-
\frac{c_v}{c_*}\Bigr)A_v+\sum_{v\in D\setminus D'}A_v ,
\]
and the last two sums are nonnegative. It remains to compute $\Sigma$. A
displacement $v\in O$ with $k$ nonzero coordinates has $c_v=p^{4-k}(p-1)^k$,
and there are $8$, $24$, $32$ and $16$ of them for $k=1,2,3,4$. The
differences $c_v-c_*$ are
\[
p^3,\qquad p^2,\qquad-p(p^2-3p+1),\qquad-2p(p-1)(p-2)+1,
\]
with the signs shown since $p\ge3$. Hence
$\Sigma=8p^3+24p^2+32p(p^2-3p+1)+16(2p(p-1)(p-2)-1)=72p^3-168p^2+96p-16$,
and with $c_0=p^4$ this is the formula for $\varepsilon$. As $-168p^2+96p
-16<0$, $\varepsilon\ge1-((1-\gamma)p+72)/(p-2)=(\gamma p-74)/(p-2)$, and if
$\gamma p\ge148$ this is at least $(\gamma p-74)/p\ge\gamma/2$.
\end{proof}

For a patch $S_s$ let $\Pi_s=\bigcup_{j\in S^\circ_s}B_j$, the set of links
with base point in $r_s+[0,md-1)^4$. For $L\ge md+2$ the patch and the
plaquettes meeting it embed in $\Z^4$, and the conditional law of
$U_{\Pi_s}$ under $\mu$ given the links outside $\Pi_s$ is the measure
$\nu^\omega_{Q}$ of Sec.~\ref{sec:massgen} for the box $Q=r_s+[0,md-1)^4$,
with $\omega$ the configuration outside. For $j\in S^\circ_s$ write
$\mathcal P^\omega_j$ for the conditional expectation under $\nu^\omega_Q$
given the links of $\Pi_s\setminus B_j$.

\begin{definition}\label{def:patchgap}
The \emph{gauge-invariant patch gap} $\gamma_{d,m}$ is the largest
$\gamma\ge0$ such that
\[
\gamma\Var_{\nu^\omega_Q}(F)\le\sum_{j\in S^\circ_s}\nu^\omega_Q\bigl((F-
\mathcal P^\omega_jF)^2\bigr)
\]
for every patch box $Q=r+[0,md-1)^4$ of $\Z^4$, every exterior
configuration $\omega$ and every $F\in\mathcal A^\omega_Q$.
\end{definition}

By translation invariance $\gamma_{d,m}$ depends only on $G$, $\varphi$,
$d$ and $m$. It is the gap, on gauge-invariant functions, of the dynamics
that resamples the $(m-1)^4$ blocks of one patch, uniform in the
configuration outside the patch.

\begin{theorem}[finite-size criterion]\label{thm:finite}
Let $d\ge2$, $m\ge4$, $d\mid L$, $L$ even, $L\ge2md$, and suppose that
$\gamma_{d,m}\ge\gamma>0$ and that the number $\varepsilon$ of
Lemma~\ref{lem:knabe} is positive. Then
\[
\frac{\varepsilon}{16}\Var_\mu(f)\le\frac1{n^4}\sum_{x\in\Lambda_L}\mu\bigl(
(f-\mathcal P_xf)^2\bigr)\qquad\text{for every real }f\in\mathcal A,
\]
that is, $\BP(2d-1,\varepsilon/16)$. If $m-1\ge148/\gamma$, then
$\varepsilon/16\ge\gamma/32$.
\end{theorem}

\begin{proof}
Fix $a$. Since $L\ge2md$, $K\ge2m\ge8$.

\emph{The patch inequality on the torus.} Let $F\in\mathcal A$; replacing $F$
by $\int F(gU)\,dg$, the integral over all gauge transformations of
$\Lambda_L$ with Haar measure, we may take $F$ invariant at every point.
Write $U=(\sigma,\omega)$ with $\sigma$ the links of $\Pi_s$ and $\omega$ the
others. Under $\mu$ the conditional law of $\sigma$ given $\omega$ is
$\nu^\omega_Q$, and for $j\in S^\circ_s$ the conditional law of $U_{B_j}$
given the other links is the same under $\mu$ and under $\nu^\omega_Q$, so
$(\mathcal P_jF)(\cdot,\omega)=\mathcal P^\omega_j(F(\cdot,\omega))$ for almost
every $\omega$ and $H_s$ acts on each fibre as $H^\omega_s=\sum_{j\in
S^\circ_s}(I-\mathcal P^\omega_j)$. A gauge transformation of $\Lambda_L$
that is trivial at the vertices not interior to $Q$ fixes $\omega$, so each
fibre $F(\cdot,\omega)$ lies in $\mathcal A^\omega_Q$. On $\mathcal
A^\omega_Q$ the operator $H^\omega_s$ is self-adjoint and nonnegative,
preserves the constants and their orthogonal complement
(Lemma~\ref{lem:boxinv}(iii)), and by Definition~\ref{def:patchgap} is at
least $\gamma$ on that complement. So its spectrum there lies in
$\{0\}\cup[\gamma,\infty)$ and $(H^\omega_s)^2\ge\gamma H^\omega_s$.
Integrating over $\omega$, $\|H_sF\|^2\ge\gamma\ip F{H_sF}$ for $F\in
\mathcal A$, that is, $H_s^2\ge\gamma H_s$ on $\mathcal A$.

\emph{Conclusion.} By Lemma~\ref{lem:knabe} with $\mathcal K=\mathcal A$,
$H^2\ge\varepsilon H$ on $\mathcal A$, so the spectrum of $H$ on $\mathcal
A$ lies in $\{0\}\cup[\varepsilon,\infty)$. If $F\in\mathcal A$ and $HF=0$,
then $\sum_j\|F-\mathcal P_jF\|^2=\ip F{HF}=0$, so $F$ is constant by
Lemma~\ref{lem:fcommute}(iv). Hence for real $f\in\mathcal A$ the function
$g=f-\mu(f)\in\mathcal A$ is orthogonal to the kernel of $H$ on $\mathcal
A$, and
\[
\varepsilon\Var_\mu(f)\le\ip g{Hg}=\sum_{j\in T_K}\mu\bigl((f-\mathcal
P_{r_j}f)^2\bigr).
\]
As $a$ runs over $[0,d)^4$ and $j$ over $T_K$, $r_j=a+dj$ runs over
$\Lambda_L$ once. Averaging over $a$ gives $\varepsilon d^4\Var_\mu(f)\le
\sum_x\mu((f-\mathcal P_xf)^2)$, and $d^4\ge n^4/16$. Since $L\ge8d\ge2n+4$
this is $\BP(2d-1,\varepsilon/16)$, and the last claim is
Lemma~\ref{lem:knabe}.
\end{proof}

The passage from the patch to the torus costs a factor: at $\varphi\equiv1$
the links are independent, $\gamma_{d,m}\ge1$ by the Efron--Stein
inequality~\cite{ES81}, and the theorem gives $\BP(2d-1,\varepsilon/16)$
with $\varepsilon<1$, against $\BP(n,1)$ from Proposition~\ref{prop:bpsup}.
